\section{Verificación Processor-in-the-Loop}

En PIL el código del controlador se compila para el procesador objetivo, se descarga a la tarjeta y
se ejecuta allí, mientras Simulink intercambia con ella las señales de entrada y salida en cada
periodo de muestreo \cite{mw_pil}. El procesador empleado fue un Arduino Uno R4 Minima, utilizado
exclusivamente como plataforma de verificación y no como el destino final del controlador.

El reparto entre lo que se ejecuta físicamente y lo que permanece simulado es el punto clave de esta
etapa. Sobre el Arduino se ejecuta únicamente el algoritmo de control: la ley PI discreta, la
saturación y el anti-windup, con el estado del integrador residiendo en la memoria del
microcontrolador. Todo lo demás sigue siendo un modelo dentro de Simulink: la planta FOPDT, el
tiempo muerto, la temperatura ambiente y la referencia. No intervienen el sensor ni el calentador de
la TCLab.

El banco \texttt{TCLab\_Test\_PIL.slx} emplea de nuevo dos ramas que referencian
\texttt{Controller.slx}, una en modo Normal y otra en modo PIL. A diferencia del banco SIL, las dos
ramas comparten solo las fuentes $\Tref$ y $\Tamb$, y cada una conserva su propia copia del retardo
y de la planta. Esa separación es necesaria porque, si las ramas compartieran la planta, cualquier
discrepancia entre las acciones de control produciría trayectorias de temperatura distintas y
dejaría de poder atribuirse la diferencia a una sola causa. Tampoco se mantuvieron simultáneamente
una instancia SIL y una PIL referenciando el mismo modelo.

Las diferencias evaluadas fueron las análogas a \eqref{eq:difsil},
$\Delta Q_1 = Q_{1,\mathrm{PIL}} - Q_{1,\mathrm{Normal}}$ y
$\Delta T_1 = T_{1,\mathrm{PIL}} - T_{1,\mathrm{Normal}}$, sobre las mismas \num{1801} muestras. El
resultado fue también de igualdad exacta (Fig.~\ref{fig:diferencias} y Tabla~\ref{tab:silpil}). No
se observó ninguna diferencia que analizar, lo cual es consistente con haber conservado el tipo
\texttt{double} y la misma estructura algorítmica en las tres ejecuciones.

\begin{figure}[!t]
\centering
\begin{tikzpicture}
\begin{groupplot}[group style={group size=1 by 2, vertical sep=0.7cm}, estilo, height=3.3cm,
                  xmin=0, xmax=1800, legend pos=north east]
\nextgroupplot[ylabel={$\Delta Q_1$ (\si{\percent})}, xticklabels={}, ymin=-1, ymax=1]
  \addplot[cAzul, line width=1.4pt, opacity=0.45] table[col sep=comma, x index=0,
    y expr=\thisrowno{2}-\thisrowno{1}] {\G/GRAFComparacionNormalSIL.csv};
  \addplot[cNaranja, dashed] table[col sep=comma, x index=0,
    y expr=\thisrowno{2}-\thisrowno{1}] {\G/GRAFComparacionNormalPIL.csv};
  \legend{SIL, PIL}
\nextgroupplot[ylabel={$\Delta T_1$ (\si{\celsius})}, xlabel={Tiempo (\si{\second})},
               ymin=-1, ymax=1]
  \addplot[cAzul, line width=1.4pt, opacity=0.45] table[col sep=comma, x index=0,
    y expr=\thisrowno{4}-\thisrowno{3}] {\G/GRAFComparacionNormalSIL.csv};
  \addplot[cNaranja, dashed] table[col sep=comma, x index=0,
    y expr=\thisrowno{4}-\thisrowno{3}] {\G/GRAFComparacionNormalPIL.csv};
  \legend{SIL, PIL}
\end{groupplot}
\end{tikzpicture}
\caption{Diferencias \emph{back-to-back} respecto de la ejecución Normal. Las cuatro trazas son
idénticamente nulas en las \num{1801} muestras; la escala vertical se fijó en
$\pm\num{1}$ para hacer visible que no hay desviación alguna.}
\label{fig:diferencias}
\end{figure}

\begin{table}[!t]
\caption{Resultados cuantitativos de las verificaciones SIL y PIL, sobre \num{1801} muestras
($T_s = \SI{1}{\second}$, \SI{1800}{\second} de simulación).}
\label{tab:silpil}
\centering
\footnotesize
\begin{tabular}{@{}lcccccc@{}}
\toprule
& \multicolumn{2}{c}{\textbf{SIL}} & \multicolumn{2}{c}{\textbf{PIL}} & \multicolumn{2}{c}{\textbf{Igualdad exacta}} \\
\cmidrule(lr){2-3}\cmidrule(lr){4-5}\cmidrule(lr){6-7}
\textbf{Señal} & Máx.\ abs. & RMSE & Máx.\ abs. & RMSE & SIL & PIL \\
\midrule
$Q_1$ (\si{\percent})   & 0 & 0 & 0 & 0 & Sí & Sí \\
$T_1$ (\si{\celsius})   & 0 & 0 & 0 & 0 & Sí & Sí \\
\midrule
\multicolumn{5}{@{}l}{\textbf{Resultado}} & \textbf{PASS} & \textbf{PASS} \\
\bottomrule
\end{tabular}
\end{table}

La evidencia que aporta PIL es distinta de la de SIL. SIL comprueba la traducción de modelo a código
usando el compilador y la aritmética del computador anfitrión; PIL añade que ese mismo código se
compila con la cadena de herramientas del procesador objetivo, cabe en su memoria y produce los
mismos resultados ejecutándose sobre él. Es decir, PIL cubre los efectos del compilador cruzado, de
los tamaños de palabra y del comportamiento aritmético del dispositivo, que SIL no puede detectar.

Aun así, un PASS en PIL no constituye una validación del controlador sobre la TCLab. Durante la
prueba la planta y el retardo siguen siendo modelos, $\Tamb$ es una señal de Simulink, no se mide
temperatura con el sensor real ni se acciona el calentador, no se ejercitan los canales físicos de
entrada y salida, y la ejecución no se realiza en tiempo real contra un proceso físico. Lo que PIL
verifica es el código y su ejecución sobre el procesador, no el lazo cerrado sobre la planta.
